\documentclass[10pt, conference]{IEEEtran}
% Nota: O BRACIS pode utilizar o formato IEEE ou Springer LNCS (Lecture Notes in Computer Science), 
% dependendo do ano ou da trilha. O padrão IEEE foi utilizado aqui como ponto de partida.

\usepackage[utf8]{inputenc}
\usepackage{cite}
\usepackage{amsmath,amssymb,amsfonts}
\usepackage{algorithmic}
\usepackage{graphicx}
\usepackage{textcomp}
\usepackage{xcolor}

\begin{document}

\title{FEDSUNT: A Novel Approach for... (Complete o Título)\\
% \thanks{Agradecimentos (opcional).}
}

\author{
\IEEEauthorblockN{1\textsuperscript{st} Nome do Autor}
\IEEEauthorblockA{\textit{Departamento} \\
\textit{Instituição}\\
Cidade, País \\
email@domain.com}
}

\maketitle

\begin{abstract}
Este artigo apresenta o FEDSUNT... (Escreva aqui o abstract do seu artigo sobre o fedsunt). 
A inteligência artificial e os sistemas inteligentes...
\end{abstract}

\begin{IEEEkeywords}
FEDSUNT, Inteligência Artificial, BRACIS, Palavra-chave 4
\end{IEEEkeywords}

\section{Introdução}
O contexto do FEDSUNT...

\section{Trabalhos Relacionados}

\section{Metodologia}

\section{Resultados e Discussão}

\section{Conclusão}

\bibliographystyle{IEEEtran}
% \bibliography{references}

\end{document}
